\section{Loss of horizontal weight} \label{sec:entropy} \label{sec:loss-hor}
In this section we prove the following purely combinatorial theorem,
which will be used to prove Theorem \ref{loss of hor weight} and ultimately Theorem \ref{def vertical teaser}. 
\begin{theorem}\label{loss-F}
Suppose $F(z) = z^2 + c$ is critically periodic of period $p$.
Let $\DD$ be the closure of the union of immediate attracting basins, and assume that $\DD$ is a union of disjoint closed Jordan domains.
Suppose that $X^n$, for $n \in \N$, are horizontal weighted arc diagrams on $\C \sm \DD$ such that
\begin{enumerate}
\item
$X^{n+1} \le X^n$, and 
\item
$F^*X^n \multimap X^{n+1}$
\end{enumerate}
for each $n \in \N$. 
Then
$$
     \| X^{8p}\|_1\leq \frac{1}{2} \| X^0 \|_1.
$$
\end{theorem}


Throughout this section,
we assume that the $\hx n$ (for $n = 0, \ldots, 8p$) are  horizontal weighted arc diagrams on $\C \sm \DD$
that satisfy Hypotheses 1 and 2 in Theorem \ref{loss-F}.

\subsection{Topological arrow}
  Let us consider two Riemann surfaces, $U\subset V$,
and let $\bal$ and $\bbe$ be multiarcs on $U$ and $V$  respectively.
We say that $\bal$ {\it topologically arrows} $\bbe$, $\bal\ta\bbe$,
if for any arc $\beta\in \bbe$
there is a sequence $(\alpha_k)$ of arcs
 with $\alpha_k \in \bal$ for each $k$,
and $(\alpha_k) \arrow \beta$.

A basic example comes from two WAD's, $X\in W(U)$ and $Y\in W(V)$,
such that $X\dominate Y$. Then  $\supp X \ta \supp Y$,
as follows immediately from the definitions.
So in the setting of this section, we have $$F^*\supp \hx {n} = \supp F^*\hx {n} \ta \supp \hx {n+1}$$ for each $n$. 

\subsection{An invariant horizontal arc diagram}


Let us say that a horizontal arc diagram $\bal$ on  $\C\sm \DD$ is {\it invariant} if
$$
   F^*\bal \ta \bal $$

\begin{prop}\label{invariance}
  There exists $n\leq 3p$ such that the horizontal arc diagram $\bal^n = \supp  \hx n$ is invariant.
\end{prop}

\begin{proof}
  Since $ \hx {n+1} \leq  \hx n$, we have
$\bal^{n+1}\subset \bal^n$.
Since $|\bal^0|\leq 3p$, there exists an  $n\leq 3p$ such that $\bal^{n+1}=\bal^n$. We take this $n$ as our desired $n$. 
Then, as we have already observed, we have $F^* \bal^n \ta \bal^{n+1} = \bal^n$.
\end{proof}

\subsection{Alignment with the Hubbard tree}
\newcommand{\bbb}{\boldsymbol \beta}
We begin with a lemma relating topological arrow (of horizontal arc diagrams) and vertical arcs. 
We observe that we can think of any vertical arc $\gamma$ in $\C \sm F^{-1}\DD$ landing at a disk in $\DD$ as a (unique and well-defined) vertical arc in $\C \sm \DD$.
\begin{lem} \label{lem:ta-int}
Suppose that 
\begin{itemize}
\item
$\bal$ is a horizontal arc diagram on $\C \sm F^{-1}\DD$,
\item
$\bbb$ is a horizontal arc diagram on $\C \sm \DD$, and 
\item
$\bal \ta \bbb$.
\end{itemize}
Suppose that $\gamma$ is a vertical arc on $\C \sm F^{-1}\DD$, landing at a disk in $\DD$, such that $\bal$ does not intersect $\gamma$.
Then $\bbb$ does not intersect $\gamma$.
\end{lem}
\begin{proof}
For every arc $\beta \in \bbb$, 
we can find a representing path $b$ such that the restriction of $b$ to $\C \sm F\inv\DD$ is a union of disjoint paths representing arcs in $\bal$. 
We can then realize $\gamma$ as a path (in $\C \sm F\inv\DD$)
 that does not intersect any of these $\bal$-representing paths. 
Therefore $\beta$ does not intersect $\gamma$.

We conclude that $\bbb$ does not intersect $\gamma$.
\end{proof}

\begin{lem}\label{H-perp}
  Any invariant horizontal arc diagram  $\bal$
is aligned with the Hubbard tree.
\end{lem}

\begin{proof}
Let  $\bal$ be an invariant horizontal arc diagram for a superattracting polynomial $F$,
and think of it for the moment as a union of actual paths,
each of which extends continuously to a map of a closed interval into $\C \sm \Int \DD$. 
Let $\De$ be the unbounded component of $\C\sm (\bal \cup \DD$).
If $\di \De$ were a subset of $\bal$,
then $\bal$ would contain a closed loop, which is absurd.
Therefore $\DD \cap \di\De$ is non-empty,
and we can find a path $b$ in $\De$ connecting $\infty$ to $\DD$. 
It represents a vertical arc $\beta_0$ such that  $\langle\bal, \beta_0\rangle=0$.

Now suppose that $\beta$ is any vertical arc for which $\langle\bal, \beta\rangle=0$,
and $\beta'$ is any lift of $\beta$ (that lands at a disk in $\DD$).
Then $\langle F^*\bal, \beta'\rangle=0$,
and, 
by Lemma \ref{lem:ta-int} and our assumption that $F^*\bal \ta \bal$,
we have $\langle \bal, \beta'\rangle=0$.
We conclude that $\pair\bal{F^*\beta} = 0$.

By induction,
we then have 
$\pair\bal{F^{*l}\beta} = 0$ for all $l$,
and hence, 
by Lemma \ref{lem:pull-to-get-Hperp},
$\pair\bal{\BH^\perp} = 0$. 
We conclude that $\bal \subset \BH$. 
\end{proof}
Putting this together with Proposition \ref{invariance} and Hypothesis 1 for the $\hx n$,
we obtain:

\begin{cor}\label{alignment}
  For any  $n\geq 3p$, the horizontal arc diagram $\bal^n = \supp  \hx n$ is
aligned with the Hubbard tree $H$.
\end{cor}


\subsection{The entropy argument}
We continue with the proof of Theorem \ref{loss-F},
where now we know that $\hx n$ is aligned with the Hubbard tree for each $n \ge 3p$. 
We first observe
\begin{lemma} \label{lem:dom-iter}
Under our standing assumption that $F^* \hx n \doms \hx {n+1}$,  we in fact have 
$F^{l*} \hx n \doms \hx {n+l}$ for all $n \in \N$ and $l \in Z^+$.
\end{lemma}
\begin{proof}
We proceed by induction,
assuming that $F^{l*} \hx n \doms \hx {n+l}$ and proving $F^{(l+1)*} \hx n \doms \hx {n+l + 1}$. 
We first observe that the assumption implies $$F^{(l+1)*} \hx n \doms F^* \hx {n+l},$$ 
and of course $$F^* \hx {n+l} \doms \hx {n+l+1}$$ by our standing assumption.
So we are finished by Lemma \ref{lem:dom-trans}.
\end{proof}

We then have, using the definition of $X|_D$ given in the preamble to Lemma \ref{lem:dom-restr},
\begin{lemma} \label{lem:res-iter}
Under our standing assumption,
$$\hx {n+l}|D_i \le (\deg F^l|D_i)\, \hx n|D_{i+l \bmod p}.$$
\end{lemma}
\begin{proof}
We have $F^{l*}\hx n|D_i = (\deg F^l|D_i)\, \hx n|D_{i+l \bmod p}$.
Moreover $F^{l*}\hx n \doms \hx {n+l}$ by Lemma \ref{lem:dom-iter}.
The Lemma then follows, by Lemma \ref{lem:dom-restr}.
\end{proof}
\begin{cor} \label{cor:sup-inf}
We have 
$$
\sup_i \hx {3n}|_{D_i} \le 2 \inf_i \hx {2n}|_{D_i} \le 2 \inf_i \hx 0|_{D_i}.
$$
\end{cor}
\begin{rem}
Lemma \ref{lem:res-iter} and its corollary is closely related to Theorem 9.3 of \cite{McM},
which relates the lengths of the hyperbolic geodesics around the small Julia sets of a given period. 
\end{rem}
We now have the following lemma, which is central to our entropy argument:
\begin{lemma} \label{lem:entropy-center}
For every $l \in \Z^+$, and $\beta \in \supp \hx {(3+l) p}$,
we have
$$
\hx {(3+l) p}(\beta) \le 2^{-l} \sup_i \hx {3p}|_{D_i}.
$$
\end{lemma}
\begin{proof}
We have $\supp X^{(3+l)p} \subseteq \BH^{lp}$ because $\supp X^{3p} \subseteq \BH$.

Suppose that $\beta$ connects $D_-$ and $D_+$.
By Lemma \ref{lem:into-flags},
there is a unique vertex-injective path in $\BH^{lp}$ connecting $D_-$ to $D_+$. 
We let the vertices of this path be $D'_0 = D_-, D'_1, \ldots, D'_n = D_+$,
and let the edges be $\gamma_0, \ldots, \gamma_{n-1}$,
where of course $D'_k \in \DD^{lp}$ and $\gamma_k \in \BH^{lp}$ for each $k$.
Moreover, 
by the same Lemma,
whenever $(\alpha_i) \to \beta$,
we can find, for each $k$, an $i$ such that $\alpha_i \in \Flag(D'_k, \gamma_k; \BH^{lp})$.
We also observe that $n \ge 2^{lp}$ by Corollary \ref{cor:shortest-lp}.

Choosing $k \in 0, \ldots, n-1$,
letting $j \in 0, \ldots, p-1$ be such that $F^{lp}(D'_k) = D_j$, and
letting $\eta$ be $F^{lp}(\gamma_k)$,
we observe that 
$$F^{lp}\from \Flag(D'_k, \gamma_k; \BH^{lp}) \to \Flag(D_j, \eta; \BH)$$
is a bijection.
Hence, 
$$
\norm{F^{lp*}\hx {3p}|_{\Flag(D'_k, \gamma_k; \BH^{lp})}}_1 = \norm{\hx {3p}|_{\Flag(D_j, \eta; \BH)}}_1
\le {\hx {3p}|_{D_i}} \le \sup_i {\hx {3p}|_{D_i}}.
$$

We then apply Lemma \ref{lem:in-each},
note the lower bound on $n$, and observe that $\bigoplus_{i=0}^{n-1} w_i \le w/n$ when $w_i \le w$ for each $i$.
\end{proof}

We can now prove Theorem \ref{loss-F}.
\begin{proof}
By Lemma \ref{lem:entropy-center} and Corollary \ref{cor:sup-inf},
letting $l = 5$, 
we have 
$$
\hx {8p}(\beta) \le  2^{-5} \sup_i \hx {3p}|_{D_i} \le 2^{-4}\inf_i \hx {0}|_{D_i}.
$$
The arc diagram $\supp \hx {8p} \cup \BH^\perp$ has at most $3p$ arcs and $\BH^\perp$ has at least $p$ arcs,
so $\hx{8p}$ has at most $2p$ arcs. 
We can then conclude
$$
\norm{\hx {8p}}_1 \le 2p \cdot 2^{-4} \inf_i \hx{0}|_{D_i} = \frac p 8 \inf_i \hx{0}|_{D_i}.
$$
On the other hand, 
$$
2\norm{\hx 0}_1 = \sum_i \hx 0|_{D_i} \ge p \inf_i \hx 0|_{D_i}. \qedhere
$$
\end{proof}

